\documentclass[10pt,aspectratio=169,serif]{beamer}

\usepackage[fontset=none, scheme=plain]{ctex}
\IfFontExistsTF{Songti SC}{%
	\setCJKmainfont{Songti SC}[AutoFakeSlant=0.2]%
	\setCJKsansfont{Songti SC}[AutoFakeSlant=0.2]%
}{%
	\setCJKmainfont{FandolSong-Regular}[Extension=.otf,BoldFont=FandolSong-Bold,AutoFakeSlant=0.2]%
	\setCJKsansfont{FandolSong-Regular}[Extension=.otf,BoldFont=FandolSong-Bold]%
}
\usepackage{amsmath,amssymb,graphicx,tikz,booktabs,multicol,calligra}
\usepackage{pgfplots}
\pgfplotsset{compat=1.18}
\usepackage{tcolorbox}
\tcbuselibrary{skins}
\usepackage{XyygU}
\usepackage{caption}
\captionsetup[figure]{name=图,labelsep=period}

\usetikzlibrary{arrows.meta,positioning,calc}

\setlength{\abovedisplayskip}{0pt plus 2pt minus 2pt}
\setlength{\belowdisplayskip}{0pt plus 2pt minus 2pt}

\title[5.5 三角恒等变换练习]{5.5 三角恒等变换（课后练习）}
\subtitle{两角和差 \& 二倍角 \& 简单恒等变换}
\author{}
\institute{}
\date{2026-10-16}

\begin{document}

\begin{frame}
  \titlepage
\end{frame}

% ================= 第一组：差角余弦公式（教材p217） =================
\begin{frame}{第一组：差角余弦公式 $C_{(\alpha-\beta)}$}
本练习对应教材第217页，巩固公式
\[
\cos(\alpha-\beta)=\cos\alpha\cos\beta+\sin\alpha\sin\beta.
\]
共5题。
\end{frame}

% ----- 第1题 -----
\begin{frame}[t]{第1题：利用 $C_{(\alpha-\beta)}$ 证明诱导公式}
\textbf{1.} 利用公式 $C_{(\alpha-\beta)}$ 证明：
\begin{enumerate}
  \item $\cos\!\left(\dfrac{3\pi}{2}-\alpha\right)=-\sin\alpha$；
  \item $\cos(-\alpha)=\cos\alpha$。
\end{enumerate}
\end{frame}

\begin{frame}[t]{第1题 解答}
\begin{tcolorbox}[colback=orange!10,colframe=orange,title=\textbf{证明}]
(1) 由 $C_{(\alpha-\beta)}$，
\[
\cos\!\left(\frac{3\pi}{2}-\alpha\right)
=\cos\frac{3\pi}{2}\cos\alpha+\sin\frac{3\pi}{2}\sin\alpha
=0\cdot\cos\alpha+(-1)\cdot\sin\alpha=-\sin\alpha.
\]
(2) $\cos(-\alpha)=\cos(0-\alpha)$，再由 $C_{(\alpha-\beta)}$：
\[
\cos(0-\alpha)=\cos 0\cos\alpha+\sin 0\sin\alpha=1\cdot\cos\alpha+0=\cos\alpha.
\]
\end{tcolorbox}
\end{frame}

% ----- 第2题 -----
\begin{frame}[t]{第2题：求 $\cos 15^\circ$}
\textbf{2.} 利用公式 $C_{(\alpha-\beta)}$ 求 $\cos 15^\circ$ 的值。
\end{frame}

\begin{frame}[t]{第2题 解答}
\begin{tcolorbox}[colback=orange!10,colframe=orange,title=\textbf{解}]
注意到 $15^\circ=45^\circ-30^\circ$，故
\[
\cos 15^\circ=\cos(45^\circ-30^\circ)
=\cos45^\circ\cos30^\circ+\sin45^\circ\sin30^\circ.
\]
代入特殊角三角函数值：
\[
=\frac{\sqrt2}{2}\cdot\frac{\sqrt3}{2}+\frac{\sqrt2}{2}\cdot\frac12
=\frac{\sqrt6+\sqrt2}{4}.
\]
\end{tcolorbox}
\end{frame}

% ----- 第3题 -----
\begin{frame}[t]{第3题：求 $\cos\!\left(\frac{\pi}{4}-\alpha\right)$}
\textbf{3.} 已知 $\cos\alpha=-\dfrac{3}{5}$，$\alpha\in\left(\dfrac{\pi}{2},\pi\right)$，求 $\cos\!\left(\dfrac{\pi}{4}-\alpha\right)$ 的值。
\end{frame}

\begin{frame}[t]{第3题 解答}
\begin{tcolorbox}[colback=orange!10,colframe=orange,title=\textbf{解}]
由 $\alpha\in(\frac{\pi}{2},\pi)$ 知 $\sin\alpha>0$：
\[
\sin\alpha=\sqrt{1-\cos^2\alpha}=\sqrt{1-\frac{9}{25}}=\frac45.
\]
由差角余弦公式：
\[
\cos\!\left(\frac{\pi}{4}-\alpha\right)
=\cos\frac{\pi}{4}\cos\alpha+\sin\frac{\pi}{4}\sin\alpha
=\frac{\sqrt2}{2}\!\left(-\frac35+\frac45\right)=\frac{\sqrt2}{10}.
\]
\end{tcolorbox}
\end{frame}

% ----- 第4题 -----
\begin{frame}[t]{第4题：求 $\cos\!\left(\theta-\frac{\pi}{3}\right)$}
\textbf{4.} 已知 $\sin\theta=\dfrac{15}{17}$，$\theta$ 是第二象限角，求 $\cos\!\left(\theta-\dfrac{\pi}{3}\right)$ 的值。
\end{frame}

\begin{frame}[t]{第4题 解答}
\begin{tcolorbox}[colback=orange!10,colframe=orange,title=\textbf{解}]
$\theta$ 为第二象限角，故 $\cos\theta<0$：
\[
\cos\theta=-\sqrt{1-\sin^2\theta}=-\sqrt{1-\frac{225}{289}}=-\frac8{17}.
\]
于是
\[
\cos\!\left(\theta-\frac{\pi}{3}\right)
=\cos\theta\cos\frac{\pi}{3}+\sin\theta\sin\frac{\pi}{3}
=-\frac8{17}\cdot\frac12+\frac{15}{17}\cdot\frac{\sqrt3}{2}
=\frac{15\sqrt3-8}{34}.
\]
\end{tcolorbox}
\end{frame}

% ----- 第5题 -----
\begin{frame}[t]{第5题：求 $\cos(\beta-\alpha)$}
\textbf{5.} 已知 $\sin\alpha=-\dfrac{2}{3}$，$\alpha\in\left(\pi,\dfrac{3\pi}{2}\right)$，$\cos\beta=\dfrac{3}{4}$，$\beta\in\left(\dfrac{3\pi}{2},2\pi\right)$，求 $\cos(\beta-\alpha)$ 的值。
\end{frame}

\begin{frame}[t]{第5题 解答}
\begin{tcolorbox}[colback=orange!10,colframe=orange,title=\textbf{解}]
由 $\alpha\in(\pi,\frac{3\pi}{2})$ 得 $\cos\alpha<0$：
\[
\cos\alpha=-\sqrt{1-\frac49}=-\frac{\sqrt5}{3}.
\]
由 $\beta\in(\frac{3\pi}{2},2\pi)$ 得 $\sin\beta<0$：
\[
\sin\beta=-\sqrt{1-\frac9{16}}=-\frac{\sqrt7}{4}.
\]
注意 $\cos(\beta-\alpha)=\cos(\alpha-\beta)$：
\[
\cos(\beta-\alpha)=\cos\alpha\cos\beta+\sin\alpha\sin\beta
=\!\left(-\frac{\sqrt5}{3}\right)\!\cdot\frac34+\!\left(-\frac23\right)\!\cdot\!\left(-\frac{\sqrt7}{4}\right)
=\frac{2\sqrt7-3\sqrt5}{12}.
\]
\end{tcolorbox}
\end{frame}

% ================= 第二组：和差角公式（教材p220） =================
\begin{frame}{第二组：和（差）角的正弦、余弦、正切公式}
本练习对应教材第220页，巩固公式
\[
\begin{aligned}
\sin(\alpha\pm\beta)&=\sin\alpha\cos\beta\pm\cos\alpha\sin\beta,\\
\cos(\alpha\pm\beta)&=\cos\alpha\cos\beta\mp\sin\alpha\sin\beta,\\
\tan(\alpha\pm\beta)&=\frac{\tan\alpha\pm\tan\beta}{1\mp\tan\alpha\tan\beta}.
\end{aligned}
\]
共5题。
\end{frame}

% ----- 第6题 -----
\begin{frame}[t]{第6题：求特殊角三角函数值}
\textbf{6.} 利用和（差）角公式，求下列各式的值：
\begin{enumerate}
  \item $\sin 15^\circ$；\qquad
  \item $\cos 75^\circ$；\qquad
  \item $\sin 75^\circ$；\qquad
  \item $\tan 15^\circ$。
\end{enumerate}
\end{frame}

\begin{frame}[t]{第6题 解答}
\begin{tcolorbox}[colback=orange!10,colframe=orange,title=\textbf{解}]
(1) $\sin15^\circ=\sin(45^\circ-30^\circ)=\sin45^\circ\cos30^\circ-\cos45^\circ\sin30^\circ
=\dfrac{\sqrt2}{2}\!\cdot\!\dfrac{\sqrt3}{2}-\dfrac{\sqrt2}{2}\!\cdot\!\dfrac12=\dfrac{\sqrt6-\sqrt2}{4}.$

(2) $\cos75^\circ=\cos(45^\circ+30^\circ)=\cos45^\circ\cos30^\circ-\sin45^\circ\sin30^\circ
=\dfrac{\sqrt6-\sqrt2}{4}.$

(3) $\sin75^\circ=\sin(45^\circ+30^\circ)=\dfrac{\sqrt6+\sqrt2}{4}.$

(4) $\tan15^\circ=\tan(45^\circ-30^\circ)=\dfrac{\tan45^\circ-\tan30^\circ}{1+\tan45^\circ\tan30^\circ}
=\dfrac{1-\frac{\sqrt3}{3}}{1+\frac{\sqrt3}{3}}=2-\sqrt3.$
\end{tcolorbox}
\end{frame}

% ----- 第7题 -----
\begin{frame}[t]{第7题：已知三角函数值求和角}
\textbf{7.}
\begin{enumerate}
  \item 已知 $\cos\theta=-\dfrac35$，$\theta\in\left(\dfrac{\pi}{2},\pi\right)$，求 $\sin\left(\theta+\dfrac{\pi}{3}\right)$ 的值；
  \item 已知 $\sin\theta=-\dfrac{12}{13}$，$\theta$ 是第三象限角，求 $\cos\left(\dfrac{\pi}{6}+\theta\right)$ 的值；
  \item 已知 $\tan\alpha=3$，求 $\tan\left(\alpha+\dfrac{\pi}{4}\right)$ 的值。
\end{enumerate}
\end{frame}

\begin{frame}[t]{第7题 解答}
\begin{tcolorbox}[colback=orange!10,colframe=orange,title=\textbf{解}]
(1) $\sin\theta=\sqrt{1-\frac{9}{25}}=\frac45$。
\[
\sin\!\left(\theta+\frac{\pi}{3}\right)=\sin\theta\cos\frac{\pi}{3}+\cos\theta\sin\frac{\pi}{3}
=\frac45\cdot\frac12+\!\left(-\frac35\right)\!\cdot\frac{\sqrt3}{2}=\frac{4-3\sqrt3}{10}.
\]
(2) $\cos\theta=-\sqrt{1-\frac{144}{169}}=-\frac5{13}$。
\[
\cos\!\left(\frac{\pi}{6}+\theta\right)=\cos\frac{\pi}{6}\cos\theta-\sin\frac{\pi}{6}\sin\theta
=\frac{\sqrt3}{2}\!\cdot\!\left(-\frac5{13}\right)-\frac12\!\cdot\!\left(-\frac{12}{13}\right)
=\frac{12-5\sqrt3}{26}.
\]
(3) $\tan(\alpha+\frac{\pi}{4})=\dfrac{\tan\alpha+\tan\frac{\pi}{4}}{1-\tan\alpha\tan\frac{\pi}{4}}=\dfrac{3+1}{1-3}=-2.$
\end{tcolorbox}
\end{frame}

% ----- 第8题 -----
\begin{frame}[t]{第8题：逆用和差公式求值}
\textbf{8.} 求下列各式的值：
\begin{multicols}{2}
\begin{enumerate}
  \item $\sin72^\circ\cos18^\circ+\cos72^\circ\sin18^\circ$；
  \item $\cos72^\circ\cos12^\circ+\sin72^\circ\sin12^\circ$；
  \item $\dfrac{\tan12^\circ+\tan33^\circ}{1-\tan12^\circ\tan33^\circ}$；
  \item $\cos74^\circ\sin14^\circ-\sin74^\circ\cos14^\circ$；
  \item $\sin34^\circ\sin26^\circ-\cos34^\circ\cos26^\circ$；
  \item $\sin20^\circ\cos110^\circ+\cos160^\circ\sin70^\circ$。
\end{enumerate}
\end{multicols}
\end{frame}

\begin{frame}[t]{第8题 解答}
\begin{tcolorbox}[colback=orange!10,colframe=orange,title=\textbf{解}]
(1) $=\sin(72^\circ+18^\circ)=\sin90^\circ=1$。

(2) $=\cos(72^\circ-12^\circ)=\cos60^\circ=\dfrac12$。

(3) $=\tan(12^\circ+33^\circ)=\tan45^\circ=1$。

(4) $=-(\sin74^\circ\cos14^\circ-\cos74^\circ\sin14^\circ)=-\sin(74^\circ-14^\circ)=-\sin60^\circ=-\dfrac{\sqrt3}{2}$。

(5) $=-(\cos34^\circ\cos26^\circ-\sin34^\circ\sin26^\circ)=-\cos(34^\circ+26^\circ)=-\cos60^\circ=-\dfrac12$。

(6) $\cos110^\circ=-\cos70^\circ,\ \cos160^\circ=-\cos20^\circ$：
原式 $=\sin20^\circ(-\cos70^\circ)+(-\cos20^\circ)\sin70^\circ=-(\sin20^\circ\cos70^\circ+\cos20^\circ\sin70^\circ)=-\sin90^\circ=-1$。
\end{tcolorbox}
\end{frame}

% ----- 第9题 -----
\begin{frame}[t]{第9题：辅助角公式化简}
\textbf{9.} 化简：
\begin{multicols}{2}
\begin{enumerate}
  \item $\dfrac12\cos x-\dfrac{\sqrt3}{2}\sin x$；
  \item $\sqrt3\sin x+\cos x$；
  \item $\sqrt2(\sin x-\cos x)$；
  \item $\sqrt2\cos x-\sqrt6\sin x$。
\end{enumerate}
\end{multicols}
\end{frame}

\begin{frame}[t]{第9题 解答}
\begin{tcolorbox}[colback=orange!10,colframe=orange,title=\textbf{解（辅助角公式）}]
(1) $\dfrac12=\cos\dfrac{\pi}{3},\ \dfrac{\sqrt3}{2}=\sin\dfrac{\pi}{3}$，故原式 $=\cos\dfrac{\pi}{3}\cos x-\sin\dfrac{\pi}{3}\sin x=\cos\!\left(x+\dfrac{\pi}{3}\right)$。

(2) 提取 $2$：$2\!\left(\dfrac{\sqrt3}{2}\sin x+\dfrac12\cos x\right)=2\sin\!\left(x+\dfrac{\pi}{6}\right)$。

(3) 提取 $\sqrt2\cdot\sqrt2=2$：$2\!\left(\dfrac{\sqrt2}{2}\sin x-\dfrac{\sqrt2}{2}\cos x\right)=2\sin\!\left(x-\dfrac{\pi}{4}\right)$。

(4) 提取 $2\sqrt2$：$2\sqrt2\!\left(\dfrac12\cos x-\dfrac{\sqrt3}{2}\sin x\right)=2\sqrt2\cos\!\left(x+\dfrac{\pi}{3}\right)$。
\end{tcolorbox}
\end{frame}

% ----- 第10题 -----
\begin{frame}[t]{第10题：求 $\sin\!\left(\beta+\frac{5\pi}{4}\right)$}
\textbf{10.} 已知 $\sin(\alpha-\beta)\cos\alpha-\cos(\beta-\alpha)\sin\alpha=\dfrac35$，$\beta$ 是第三象限角，求 $\sin\!\left(\beta+\dfrac{5\pi}{4}\right)$ 的值。
\end{frame}

\begin{frame}[t]{第10题 解答}
\begin{tcolorbox}[colback=orange!10,colframe=orange,title=\textbf{解}]
注意 $\cos(\beta-\alpha)=\cos(\alpha-\beta)$，故左端
\[
\sin(\alpha-\beta)\cos\alpha-\cos(\alpha-\beta)\sin\alpha=\sin[(\alpha-\beta)-\alpha]=\sin(-\beta)=-\sin\beta.
\]
由 $-\sin\beta=\frac35$ 得 $\sin\beta=-\frac35$。$\beta$ 在第三象限，$\cos\beta=-\frac45$。
\[
\sin\!\left(\beta+\frac{5\pi}{4}\right)=\sin\beta\cos\frac{5\pi}{4}+\cos\beta\sin\frac{5\pi}{4}
=\!\left(-\frac35\right)\!\cdot\!\left(-\frac{\sqrt2}{2}\right)+\!\left(-\frac45\right)\!\cdot\!\left(-\frac{\sqrt2}{2}\right)
=\frac{7\sqrt2}{10}.
\]
\end{tcolorbox}
\end{frame}

% ================= 第三组：二倍角公式（教材p223） =================
\begin{frame}{第三组：二倍角的正弦、余弦、正切公式}
本练习对应教材第223页，巩固
\[
\sin 2\alpha=2\sin\alpha\cos\alpha,\quad
\cos 2\alpha=\cos^2\alpha-\sin^2\alpha=1-2\sin^2\alpha=2\cos^2\alpha-1,\quad
\tan 2\alpha=\frac{2\tan\alpha}{1-\tan^2\alpha}.
\]
共5题。
\end{frame}

% ----- 第11题 -----
\begin{frame}[t]{第11题：半角求倍角}
\textbf{11.} 已知 $\cos\dfrac{\alpha}{8}=-\dfrac45$，$8\pi<\alpha<12\pi$，求 $\sin\dfrac{\alpha}{4}$，$\cos\dfrac{\alpha}{4}$，$\tan\dfrac{\alpha}{4}$ 的值。
\end{frame}

\begin{frame}[t]{第11题 解答}
\begin{tcolorbox}[colback=orange!10,colframe=orange,title=\textbf{解}]
由 $8\pi<\alpha<12\pi$，得 $4\pi<\dfrac{\alpha}{2}<6\pi$，$2\pi<\dfrac{\alpha}{4}<3\pi$，$\pi<\dfrac{\alpha}{8}<\dfrac{3\pi}{2}$。

可见 $\dfrac{\alpha}{8}$ 在第三象限（与 $\cos\dfrac{\alpha}{8}=-\dfrac45$ 一致），$\dfrac{\alpha}{4}$ 在第三象限（$2\pi$ 到 $3\pi$ 即 $\pi$ 到 $2\pi$ 加 $2\pi$），故 $\sin\dfrac{\alpha}{4}<0,\ \cos\dfrac{\alpha}{4}<0$。

由二倍角公式：
\[
\sin\frac{\alpha}{4}=2\sin\frac{\alpha}{8}\cos\frac{\alpha}{8}.
\]
$\dfrac{\alpha}{8}$ 在第三象限，$\sin\dfrac{\alpha}{8}=-\sqrt{1-\frac{16}{25}}=-\dfrac35$。
\[
\sin\frac{\alpha}{4}=2\cdot\!\left(-\frac35\right)\!\cdot\!\left(-\frac45\right)=\frac{24}{25},\quad
\cos\frac{\alpha}{4}=2\cos^2\frac{\alpha}{8}-1=2\cdot\frac{16}{25}-1=\frac7{25}.
\]
$\tan\dfrac{\alpha}{4}=\dfrac{\sin(\alpha/4)}{\cos(\alpha/4)}=\dfrac{24}{7}$。
\end{tcolorbox}
\end{frame}

% ----- 第12题 -----
\begin{frame}[t]{第12题：由 $\sin(\alpha-\pi)$ 求 $\cos 2\alpha$}
\textbf{12.} 已知 $\sin(\alpha-\pi)=\dfrac35$，求 $\cos 2\alpha$ 的值。
\end{frame}

\begin{frame}[t]{第12题 解答}
\begin{tcolorbox}[colback=orange!10,colframe=orange,title=\textbf{解}]
由诱导公式 $\sin(\alpha-\pi)=-\sin(\alpha-\pi+\pi)=-\sin\alpha$，故
\[
-\sin\alpha=\frac35\ \Longrightarrow\ \sin\alpha=-\frac35.
\]
于是
\[
\cos 2\alpha=1-2\sin^2\alpha=1-2\cdot\frac{9}{25}=\frac{7}{25}.
\]
\end{tcolorbox}
\end{frame}

% ----- 第13题 -----
\begin{frame}[t]{第13题：由 $\sin 2\alpha=-\sin\alpha$ 求 $\tan\alpha$}
\textbf{13.} 已知 $\sin 2\alpha=-\sin\alpha$，$\alpha\in\left(\dfrac{\pi}{2},\pi\right)$，求 $\tan\alpha$ 的值。
\end{frame}

\begin{frame}[t]{第13题 解答}
\begin{tcolorbox}[colback=orange!10,colframe=orange,title=\textbf{解}]
$\sin 2\alpha=2\sin\alpha\cos\alpha=-\sin\alpha$。
因 $\alpha\in(\frac{\pi}{2},\pi)$，$\sin\alpha\neq 0$，两边除以 $\sin\alpha$：
\[
2\cos\alpha=-1\ \Longrightarrow\ \cos\alpha=-\frac12.
\]
$\alpha\in(\frac{\pi}{2},\pi)$，故 $\alpha=\frac{2\pi}{3}$，$\sin\alpha=\frac{\sqrt3}{2}$。
\[
\tan\alpha=\frac{\sin\alpha}{\cos\alpha}=\frac{\sqrt3/2}{-1/2}=-\sqrt3.
\]
\end{tcolorbox}
\end{frame}

% ----- 第14题 -----
\begin{frame}[t]{第14题：由 $\tan 2\alpha=\frac13$ 求 $\tan\alpha$}
\textbf{14.} 已知 $\tan 2\alpha=\dfrac13$，求 $\tan\alpha$ 的值。
\end{frame}

\begin{frame}[t]{第14题 解答}
\begin{tcolorbox}[colback=orange!10,colframe=orange,title=\textbf{解}]
由二倍角正切公式 $\tan 2\alpha=\dfrac{2\tan\alpha}{1-\tan^2\alpha}$，令 $t=\tan\alpha$：
\[
\frac{2t}{1-t^2}=\frac13\ \Longrightarrow\ 6t=1-t^2\ \Longrightarrow\ t^2+6t-1=0.
\]
解得 $t=-3\pm\sqrt{10}$。
\[
\alert{\tan\alpha=-3+\sqrt{10}\quad\text{或}\quad\tan\alpha=-3-\sqrt{10}}.
\]
（两解均合理，对应 $2\alpha$ 在不同象限。）
\end{tcolorbox}
\end{frame}

% ----- 第15题 -----
\begin{frame}[t]{第15题：逆用二倍角公式求值}
\textbf{15.} 求下列各式的值：
\begin{enumerate}
  \item $\sin15^\circ\cos15^\circ$；\qquad
  \item $\cos^2\dfrac{\pi}{8}-\sin^2\dfrac{\pi}{8}$；
  \item $\dfrac{\tan22.5^\circ}{1-\tan^2 22.5^\circ}$；\qquad
  \item $2\cos^2 22.5^\circ-1$。
\end{enumerate}
\end{frame}

\begin{frame}[t]{第15题 解答}
\begin{tcolorbox}[colback=orange!10,colframe=orange,title=\textbf{解}]
(1) $\sin15^\circ\cos15^\circ=\dfrac12\sin30^\circ=\dfrac14$。

(2) $\cos^2\dfrac{\pi}{8}-\sin^2\dfrac{\pi}{8}=\cos\dfrac{\pi}{4}=\dfrac{\sqrt2}{2}$。

(3) $\dfrac{\tan22.5^\circ}{1-\tan^2 22.5^\circ}=\dfrac12\tan45^\circ=\dfrac12$。

(4) $2\cos^2 22.5^\circ-1=\cos45^\circ=\dfrac{\sqrt2}{2}$。
\end{tcolorbox}
\end{frame}

% ================= 第四组：简单恒等变换（教材p226） =================
\begin{frame}{第四组：半角公式、积化和差、和差化积}
本练习对应教材第226页，巩固
\[
\sin^2\frac{\alpha}{2}=\frac{1-\cos\alpha}{2},\quad
\cos^2\frac{\alpha}{2}=\frac{1+\cos\alpha}{2},\quad
\tan\frac{\alpha}{2}=\frac{\sin\alpha}{1+\cos\alpha}=\frac{1-\cos\alpha}{\sin\alpha}.
\]
共5题。
\end{frame}

% ----- 第16题 -----
\begin{frame}[t]{第16题：半角正切公式证明}
\textbf{16.} 求证：
\[
\tan\frac{\alpha}{2}=\frac{\sin\alpha}{1+\cos\alpha}=\frac{1-\cos\alpha}{\sin\alpha}.
\]
\end{frame}

\begin{frame}[t]{第16题 解答}
\begin{tcolorbox}[colback=orange!10,colframe=orange,title=\textbf{证明}]
利用二倍角公式 $\sin\alpha=2\sin\dfrac{\alpha}{2}\cos\dfrac{\alpha}{2}$，$1+\cos\alpha=2\cos^2\dfrac{\alpha}{2}$，$1-\cos\alpha=2\sin^2\dfrac{\alpha}{2}$：
\[
\frac{\sin\alpha}{1+\cos\alpha}=\frac{2\sin\frac{\alpha}{2}\cos\frac{\alpha}{2}}{2\cos^2\frac{\alpha}{2}}=\frac{\sin\frac{\alpha}{2}}{\cos\frac{\alpha}{2}}=\tan\frac{\alpha}{2}.
\]
\[
\frac{1-\cos\alpha}{\sin\alpha}=\frac{2\sin^2\frac{\alpha}{2}}{2\sin\frac{\alpha}{2}\cos\frac{\alpha}{2}}=\frac{\sin\frac{\alpha}{2}}{\cos\frac{\alpha}{2}}=\tan\frac{\alpha}{2}.
\]
\end{tcolorbox}
\end{frame}

% ----- 第17题 -----
\begin{frame}[t]{第17题：已知 $\cos\theta$ 求半角}
\textbf{17.} 已知 $\cos\theta=\dfrac13$，且 $270^\circ<\theta<360^\circ$，试求 $\sin\dfrac{\theta}{2}$ 和 $\cos\dfrac{\theta}{2}$ 的值。
\end{frame}

\begin{frame}[t]{第17题 解答}
\begin{tcolorbox}[colback=orange!10,colframe=orange,title=\textbf{解}]
由 $270^\circ<\theta<360^\circ$，得 $135^\circ<\dfrac{\theta}{2}<180^\circ$，即 $\dfrac{\theta}{2}$ 在第二象限，$\sin\dfrac{\theta}{2}>0$，$\cos\dfrac{\theta}{2}<0$。
\[
\sin\frac{\theta}{2}=\sqrt{\frac{1-\cos\theta}{2}}=\sqrt{\frac{1-\frac13}{2}}=\sqrt{\frac13}=\frac{\sqrt3}{3}.
\]
\[
\cos\frac{\theta}{2}=-\sqrt{\frac{1+\cos\theta}{2}}=-\sqrt{\frac{1+\frac13}{2}}=-\sqrt{\frac23}=-\frac{\sqrt6}{3}.
\]
\end{tcolorbox}
\end{frame}

% ----- 第18题 -----
\begin{frame}[t]{第18题：等腰三角形底角正切}
\textbf{18.} 已知等腰三角形的顶角的余弦等于 $\dfrac{7}{25}$，求这个三角形的一个底角的正切。
\end{frame}

\begin{frame}[t]{第18题 解答}
\begin{tcolorbox}[colback=orange!10,colframe=orange,title=\textbf{解}]
设顶角为 $\alpha$，底角为 $\beta$，则 $\alpha+2\beta=\pi$，$\beta=\dfrac{\pi-\alpha}{2}=\dfrac{\pi}{2}-\dfrac{\alpha}{2}$。
\[
\tan\beta=\tan\!\left(\frac{\pi}{2}-\frac{\alpha}{2}\right)=\cot\frac{\alpha}{2}=\frac{1}{\tan(\alpha/2)}.
\]
由半角公式 $\tan\dfrac{\alpha}{2}=\sqrt{\dfrac{1-\cos\alpha}{1+\cos\alpha}}=\sqrt{\dfrac{1-7/25}{1+7/25}}=\sqrt{\dfrac{18}{32}}=\dfrac34$。
故 $\tan\beta=\dfrac43$。
\end{tcolorbox}
\end{frame}

% ----- 第19题 -----
\begin{frame}[t]{第19题：积化和差公式证明}
\textbf{19.} 求证：
\begin{enumerate}
  \item $\cos\alpha\sin\beta=\dfrac12[\sin(\alpha+\beta)-\sin(\alpha-\beta)]$；
  \item $\cos\alpha\cos\beta=\dfrac12[\cos(\alpha+\beta)+\cos(\alpha-\beta)]$；
  \item $\sin\alpha\sin\beta=-\dfrac12[\cos(\alpha+\beta)-\cos(\alpha-\beta)]$。
\end{enumerate}
\end{frame}

\begin{frame}[t]{第19题 解答}
\begin{tcolorbox}[colback=orange!10,colframe=orange,title=\textbf{证明}]
由 $\sin(\alpha+\beta)=\sin\alpha\cos\beta+\cos\alpha\sin\beta$ 与 $\sin(\alpha-\beta)=\sin\alpha\cos\beta-\cos\alpha\sin\beta$，两式相减：
\[
\sin(\alpha+\beta)-\sin(\alpha-\beta)=2\cos\alpha\sin\beta,
\]
即得 (1)。

由 $\cos(\alpha+\beta)=\cos\alpha\cos\beta-\sin\alpha\sin\beta$，$\cos(\alpha-\beta)=\cos\alpha\cos\beta+\sin\alpha\sin\beta$，两式相加得 (2)；两式相减得
\[
\cos(\alpha-\beta)-\cos(\alpha+\beta)=2\sin\alpha\sin\beta,
\]
即得 (3)。
\end{tcolorbox}
\end{frame}

% ----- 第20题 -----
\begin{frame}[t]{第20题：和差化积公式证明}
\textbf{20.} 求证：
\begin{enumerate}
  \item $\sin\theta-\sin\varphi=2\cos\dfrac{\theta+\varphi}{2}\sin\dfrac{\theta-\varphi}{2}$；
  \item $\cos\theta+\cos\varphi=2\cos\dfrac{\theta+\varphi}{2}\cos\dfrac{\theta-\varphi}{2}$；
  \item $\cos\theta-\cos\varphi=-2\sin\dfrac{\theta+\varphi}{2}\sin\dfrac{\theta-\varphi}{2}$。
\end{enumerate}
\end{frame}

\begin{frame}[t]{第20题 解答}
\begin{tcolorbox}[colback=orange!10,colframe=orange,title=\textbf{证明}]
令 $\alpha=\dfrac{\theta+\varphi}{2}$，$\beta=\dfrac{\theta-\varphi}{2}$，则 $\theta=\alpha+\beta$，$\varphi=\alpha-\beta$。

由 $\sin(\alpha+\beta)-\sin(\alpha-\beta)=2\cos\alpha\sin\beta$，代入即得 (1)。

由 $\cos(\alpha+\beta)+\cos(\alpha-\beta)=2\cos\alpha\cos\beta$，代入即得 (2)。

由 $\cos(\alpha+\beta)-\cos(\alpha-\beta)=-2\sin\alpha\sin\beta$，代入即得 (3)。
\end{tcolorbox}
\end{frame}

% ================= 第五组：辅助角应用（教材p228） =================
\begin{frame}{第五组：辅助角公式的应用（教材p228）}
本练习对应教材第228页，巩固 $a\sin x+b\cos x=\sqrt{a^2+b^2}\sin(x+\varphi)$ 化一角一函数。
共3题。
\end{frame}

% ----- 第21题 -----
\begin{frame}[t]{第21题：求周期、最大值、最小值}
\textbf{21.} 求下列函数的周期、最大值和最小值：
\begin{enumerate}
  \item $y=5\cos x-12\sin x$；
  \item $y=\cos x+2\sin x$。
\end{enumerate}
\end{frame}

\begin{frame}[t]{第21题 解答}
\begin{tcolorbox}[colback=orange!10,colframe=orange,title=\textbf{解}]
(1) $A=\sqrt{5^2+12^2}=13$，故 $y=13\cos(x+\varphi)$（其中 $\cos\varphi=\frac5{13},\sin\varphi=\frac{12}{13}$）。
周期 $T=2\pi$，最大值 $13$，最小值 $-13$。

(2) $A=\sqrt{1^2+2^2}=\sqrt5$，故 $y=\sqrt5\sin(x+\varphi)$。
周期 $T=2\pi$，最大值 $\sqrt5$，最小值 $-\sqrt5$。
\end{tcolorbox}
\end{frame}

% ----- 第22题 -----
\begin{frame}[t]{第22题：圆内接矩形最大面积}
\textbf{22.} 要在半径为 $R$ 的圆形场地内建一个矩形的花坛，应怎样截取，才能使花坛的面积最大？
\end{frame}

\begin{frame}[t]{第22题 解答}
\begin{tcolorbox}[colback=orange!10,colframe=orange,title=\textbf{解}]
圆内接矩形的对角线为圆的直径 $2R$。设矩形一边与直径夹角为 $\theta$，则
\[
a=2R\cos\theta,\qquad b=2R\sin\theta.
\]
面积
\[
S=ab=4R^2\sin\theta\cos\theta=2R^2\sin 2\theta.
\]
当 $\sin 2\theta=1$ 即 $\theta=\dfrac{\pi}{4}$ 时，$S_{\max}=2R^2$。
此时 $a=b=\sqrt2 R$，即花坛取圆内接正方形时面积最大，最大面积为 $2R^2$。
\end{tcolorbox}
\end{frame}

% ----- 第23题 -----
\begin{frame}[t]{第23题：正 $n$ 边形的外接圆与内切圆}
\textbf{23.} 已知正 $n$ 边形的边长为 $a$，内切圆半径为 $r$，外接圆半径为 $R$。求证：
\[
R+r=\frac{a}{2\tan\frac{\pi}{2n}}.
\]
\end{frame}

\begin{frame}[t]{第23题 解答}
\begin{tcolorbox}[colback=orange!10,colframe=orange,title=\textbf{证明}]
设正 $n$ 边形的中心角为 $\dfrac{2\pi}{n}$。由中心向一边作垂线（即内切圆半径 $r$），该边被平分，半边长为 $\dfrac{a}{2}$。在直角三角形中：
\[
\frac{a/2}{r}=\tan\frac{\pi}{n},\quad \frac{a/2}{R}=\sin\frac{\pi}{n}.
\]
故
\[
r=\frac{a}{2\tan\frac{\pi}{n}},\qquad R=\frac{a}{2\sin\frac{\pi}{n}}.
\]
于是
\[
R+r=\frac{a}{2}\!\left(\frac{1}{\sin\frac{\pi}{n}}+\frac{\cos\frac{\pi}{n}}{\sin\frac{\pi}{n}}\right)
=\frac{a}{2}\cdot\frac{1+\cos\frac{\pi}{n}}{\sin\frac{\pi}{n}}.
\]
由半角公式 $\dfrac{1+\cos\alpha}{\sin\alpha}=\cot\dfrac{\alpha}{2}$，取 $\alpha=\dfrac{\pi}{n}$，得
\[
R+r=\frac{a}{2}\cot\frac{\pi}{2n}=\frac{a}{2\tan\frac{\pi}{2n}}.
\]
\end{tcolorbox}
\end{frame}

\begin{frame}[c]
  \begin{center}
    {\Huge\calligra Thanks!}
  \end{center}
\end{frame}

\end{document}
